% capitulo1.tex

\tituloI{EL PROBLEMA DE INVESTIGACIÓN}

\tituloII{Planteamiento del Problema}

\parrafoII{
A nivel global, la integración de la Inteligencia Artificial Generativa (IAG) en el ciclo de vida del desarrollo de software ha transformado la productividad de la industria; no obstante, ha introducido un fenómeno crítico denominado "deuda técnica automatizada". Investigaciones recientes como la de Hou et al. \cite{hou2024large} sostienen que el código generado por modelos de lenguaje de gran escala (LLMs) a menudo carece de una visión arquitectónica integral, derivando en sistemas con una complejidad estructural elevada. Esta tendencia sugiere un incremento sustancial en la velocidad de codificación, pero a costa de una degradación en la mantenibilidad, lo que contraviene los estándares de calidad interna establecidos por la norma ISO/IEC 25010 \cite{iso25010}.
}

\parrafoII{
En el contexto nacional, el Perú atraviesa un proceso de aceleración digital orientado a la optimización de procesos mediante la automatización. Según el \cite{inei2024tecnologia}, existe un incremento sustancial en la adopción de servicios tecnológicos en el sector empresarial; sin embargo, este avance presenta desafíos estructurales. De acuerdo con la \cite{sgtd2024estrategia}, el país enfrenta una brecha en la formación de talento especializado y en la definición de estándares de gobernanza técnica que aseguren la calidad de las soluciones de IA. Esta carencia de marcos de validación estandarizados eleva el riesgo de acumular deuda técnica en sistemas críticos, comprometiendo la mantenibilidad y la seguridad del software desarrollado en el entorno nacional.
}

\parrafoII{
A nivel local, se observa en el entorno académico y profesional de la región Cusco una adopción empírica de la Inteligencia Artificial Generativa para la resolución de problemas de programación. A través de la experiencia en el desarrollo de proyectos de software dentro de la Universidad Continental, se ha identificado que la prioridad suele centrarse en la funcionalidad inmediata, postergando la calidad interna del código. Esta observación coincide con las brechas identificadas en el \cite{gore2025plan}, donde se evidencia que la transformación digital de la región requiere de una mayor rigurosidad en los estándares técnicos. Por lo tanto, la presente investigación surge de la necesidad de establecer herramientas que transformen esta práctica empírica observada en la formación profesional en un proceso estandarizado que asegure la mantenibilidad según la norma ISO/IEC 25010.
}


\parrafoII{
El problema central se identifica como la insuficiente calidad estructural y la elevada deuda técnica presente en el software generado mediante Inteligencia Artificial Generativa (IAG) ante la ausencia de marcos metodológicos de ingeniería. Investigaciones previas \cite{liu2023refined} demuestran que, aunque los modelos de lenguaje producen código funcional, este suele presentar deficiencias en los atributos de mantenibilidad y complejidad ciclomática definidos por la norma ISO/IEC 25010. Asimismo, la carencia de un método sistemático para estructurar los requerimientos técnicos o prompts deriva en lo que la literatura denomina como opacidad técnica o ``cajas negras'' de código \cite{siddiq2023exploring}, dificultando los procesos de auditoría, optimización y escalabilidad del software en entornos de producción.
}

\parrafoII{
Las causas de esta problemática son multidimensionales. En el plano técnico, se observa una dependencia de modelos de IA de acceso abierto que, si bien son potentes, requieren de una guía estructural para evitar la redundancia lógica. A nivel organizacional, existe una presión por reducir los tiempos de entrega (Time-to-Market), lo que lleva a omitir las fases de análisis estático y validación de arquitectura. En el factor humano, destaca la carencia de competencias en ingeniería de prompts aplicada a la ingeniería de software, tratando a la IA como una herramienta de consulta y no como un componente de un pipeline de producción sistémico.
}

\parrafoII{
Las consecuencias de persistir en este modelo empírico de desarrollo son graves. Operativamente, las organizaciones enfrentarán un incremento exponencial en los costos de mantenimiento preventivo y correctivo. Económicamente, la deuda técnica acumulada puede superar el valor del activo de software, llevando a la quiebra técnica de proyectos. Socialmente, la baja confiabilidad de los sistemas puede resultar en la pérdida de integridad de datos sensibles y una desconfianza generalizada hacia la digitalización basada en IA. No se trata solo de un error de sintaxis, sino de un compromiso de la resiliencia organizacional.
}


\parrafoII{
De no establecerse un **Framework de Ingeniería de Prompts** que actúe como un filtro de calidad basado en la norma \cite{iso25010}, la industria del software se enfrentará a una saturación de sistemas heredados (legacy) generados por IA que serán imposibles de mantener. Prospectivamente, el éxito de la ingeniería de sistemas en la próxima década dependerá de la capacidad de los profesionales para sistematizar la interacción con la IA, transformando la generación de código de un proceso aleatorio a una disciplina de ingeniería predictiva y controlada, asegurando que la tecnología sea un motor de eficiencia y no un generador de inestabilidad sistémica.
}
% capitulo1_formulacion.tex

\tituloII{Formulación del Problema}

\tituloIII{Problema General}

\parrafoII{
¿Qué efecto tiene la implementación de un Framework de Ingeniería de Prompts basado en la norma ISO/IEC 25010 en la mantenibilidad y la deuda técnica del código generado por inteligencia artificial generativa?
}

\tituloIII{Problemas Específicos}

\begin{listaPE}
    \item  ¿Qué efecto tiene la aplicación de patrones arquitectónicos de prompts en la complejidad ciclomática de McCabe del código generado?
    
    \item  ¿Qué efecto tiene la implementación de un pipeline de análisis estático en la reducción de la deuda técnica del software generado?
    
    \item  ¿Qué efecto tiene la aplicación del protocolo de validación basado en ISO/IEC 25010 en el índice de mantenibilidad del código generado?
\end{listaPE}

% capitulo1_objetivos.tex

\tituloII{Objetivos de la Investigación}

\tituloIII{Objetivo General}

\parrafoIII{
Evaluar el efecto de la implementación de un Framework de Ingeniería de Prompts basado en la norma ISO/IEC 25010 en la mantenibilidad y la deuda técnica del código generado por inteligencia artificial generativa.
}

\tituloIII{Objetivos Específicos}

\begin{listaOE}
    \item Evaluar el efecto de la aplicación de patrones arquitectónicos de prompts en la complejidad ciclomática de McCabe.
    
    \item Evaluar el efecto de la implementación de un pipeline de análisis estático en la reducción de la deuda técnica.
    
    \item Evaluar el efecto de la aplicación del protocolo de validación basado en ISO/IEC 25010 en el índice de mantenibilidad.
\end{listaOE}

% capitulo1_justificacion.tex

\tituloII{Justificación de la Investigación}

\tituloIII{Justificación Teórica}

\parrafoIII{
Desde una perspectiva teórica, esta investigación se justifica por la necesidad de integrar los paradigmas de la Inteligencia Artificial Generativa con los modelos clásicos de calidad de software. El estudio profundiza en la relación entre la estructuración de instrucciones en lenguaje natural (Ingeniería de Prompts) y la complejidad estructural del código fuente, utilizando como base la teoría de la Complejidad Ciclomática de McCabe y los pilares de mantenibilidad de la norma ISO/IEC 25010. Este trabajo aporta un nuevo marco conceptual que permite entender a la IA no solo como una herramienta de productividad, sino como un componente sistémico que requiere de rigor arquitectónico para garantizar la sostenibilidad tecnológica a largo plazo.
}

\tituloIII{Justificación Práctica}

\parrafoIII{
En el ámbito práctico, la investigación resuelve el problema crítico de la acumulación de deuda técnica en los departamentos de TI que han adoptado el uso empírico de la IA. Al proporcionar un Framework de Ingeniería de Prompts validado, se dota a los ingenieros de software de una herramienta metodológica que reduce los tiempos de refactorización y los costos operativos asociados al mantenimiento preventivo. La implementación de este marco permite que las soluciones de software generadas sean profesionalmente aptas para entornos de producción, minimizando los riesgos de fallos lógicos y optimizando el ciclo de vida de desarrollo en las organizaciones tecnológicas actuales.
}

\tituloIII{Justificación Metodológica}

\parrafoIII{
Metodológicamente, el estudio propone un protocolo de experimentación y medición sistemática que puede ser replicado en futuros estudios sobre ingeniería de software asistida. Se justifica mediante la creación de un \textit{pipeline} de evaluación que combina el análisis estático de código con la validación normativa internacional. Esta propuesta establece procedimientos claros para cuantificar la calidad del código generado, transformando una actividad que actualmente es subjetiva y variable en un proceso de ingeniería controlado, medible y auditable, sentando un precedente para la estandarización del uso de LLMs en la formación académica y profesional.
}
% capitulo1_delimitacion.tex

\tituloII{Delimitación de la Investigación}

\tituloIII{Delimitación Espacial}

\parrafoIII{
La investigación se llevará a cabo en el Laboratorio de Ingeniería de Sistemas de la Universidad Continental, sede Cusco. El entorno de experimentación estará constituido por nodos computacionales locales (clústeres de pruebas) bajo arquitectura Linux, donde se desplegarán los modelos de lenguaje y las herramientas de análisis estático de código para garantizar un control estricto sobre las variables de red y procesamiento.
}

\tituloIII{Delimitación Temporal}

\parrafoIII{
El estudio se desarrollará durante el periodo académico 2026-I, comprendido entre los meses de marzo y julio de 2026. Este cronograma incluye las fases de diseño del Framework, la ejecución de las pruebas experimentales automatizadas, el procesamiento de métricas de calidad y la redacción del informe final de resultados.
}

\tituloIII{Delimitación Social}

\parrafoIII{
La investigación está dirigida a la comunidad de ingeniería de software, específicamente a desarrolladores, arquitectos de sistemas y analistas de calidad (QA). Los resultados beneficiarán directamente a las organizaciones tecnológicas locales en Cusco que buscan estandarizar el uso de herramientas de inteligencia artificial generativa bajo marcos de cumplimiento internacional.
}

\tituloII{Hipótesis y Variables}

\tituloIII{Hipótesis}

\parrafoIII{\textbf{Hipótesis General:} La implementación de un Framework de Ingeniería de Prompts basado en la norma ISO/IEC 25010 reduce significativamente la deuda técnica y mejora la mantenibilidad del código generado por inteligencia artificial generativa en comparación con métodos empíricos.}

\parrafoIII{\textbf{Hipótesis Específicas:}}
\begin{listaHE}
    \item El código generado mediante el uso de patrones arquitectónicos de prompts presenta menor complejidad ciclomática que el código generado mediante métodos empíricos.
    \item El código generado con un pipeline de análisis estático presenta menor deuda técnica que el generado sin dicho pipeline.
    \item El código generado bajo el protocolo de validación basado en ISO/IEC 25010 presenta un mayor índice de mantenibilidad que el código generado sin dicho protocolo.
\end{listaHE}

\tituloIII{Definición de Variables}

\parrafoIII{\textbf{Variable Independiente (X):} Aplicación del Framework de Ingeniería de Prompts}
\begin{itemize}[nosep]
        \item Nivel 1: Sin framework (método empírico)
        \item Nivel 2: Con framework estructurado
    \end{itemize}

\parrafoIII{\textbf{Variables Dependientes (Y):}}
\begin{listaVD}
    \item \textbf{Mantenibilidad del Software:} Medida a través del índice de mantenibilidad y la complejidad ciclomática de McCabe utilizando herramientas de análisis estático (SonarQube/Radon).
    \item \textbf{Deuda Técnica:} Medida a través del tiempo estimado de refactorización y la densidad de defectos detectados por el \textit{pipeline} automatizado.
    %\item \textbf{Fiabilidad Sistémica:} Medida a través del porcentaje de cumplimiento de los requisitos de idoneidad técnica establecidos en la norma ISO/IEC 25010.
\end{listaVD}